% ---------------------------------------------------------------------------
% preamble.tex -- shared packages, styles and macros for the LuPNT manual.
%
% Every chapter under chapters/ is written against the contract defined here:
% the listing styles, the \code/\file/\cpp/\py macros, and the label prefixes
% (chap:, sec:, eq:, tab:, fig:).  Chapters must not load packages of their
% own -- add them here so that a chapter stays compilable in isolation via
% \includeonly.
% ---------------------------------------------------------------------------

\usepackage[T1]{fontenc}
\usepackage[utf8]{inputenc}
\usepackage{lmodern}
\usepackage[margin=1in]{geometry}

% --- mathematics ------------------------------------------------------------
\usepackage{amsmath}
\usepackage{amssymb}
\usepackage{bm}
\usepackage{siunitx}
% \year and \day are TeX count registers, so these units need their own
% control sequences rather than the siunitx defaults.
% siunitx's default \micro resolves to \textmu, which under T1 + Latin Modern
% renders as a t-cedilla ("0.3 tps") rather than a mu. Every microsecond,
% micrometre and microradian in the manual goes through this, so redeclare the
% prefix to use the math mu. Verified by rendering, not by text extraction:
% pdftotext maps the broken glyph misleadingly either way.
\DeclareSIPrefix{\micro}{\ensuremath{\mu}}{-6}

\DeclareSIUnit{\yr}{yr}
\DeclareSIUnit{\dayunit}{d}
\DeclareSIUnit{\bit}{bit}
\DeclareSIUnit{\dB}{dB}
\DeclareSIUnit{\dBHz}{dB\text{-}Hz}
\DeclareSIUnit{\dBW}{dBW}
\DeclareSIUnit{\dBi}{dBi}
\DeclareSIUnit{\TECU}{TECU}
\DeclareSIUnit{\au}{au}

% --- tables, figures, floats ------------------------------------------------
\usepackage{booktabs}
\usepackage{longtable}
\usepackage{array}
\usepackage{multirow}
\usepackage{tabularx}
\usepackage{graphicx}
\usepackage{float}

% --- where figures come from ------------------------------------------------
% latex_docs holds text only. Binary artefacts live in one of two places:
%
%   ../assets/               repository-level brand assets, already in git
%   $LUPNT_DATA_PATH/docs/   generated/large artefacts, fetched not committed
%
% build.sh resolves LUPNT_DATA_PATH and writes figpath.tex (gitignored) with
% the absolute path; the \providecommand below is the fallback for a plain
% `pdflatex lupnt` run from this directory.
\IfFileExists{figpath.tex}{\input{figpath}}{}
\providecommand{\datafigdir}{../data/LuPNT_data/docs/}
\graphicspath{{\datafigdir}{../assets/}{../docs/_static/}}

% \datafig[<includegraphics opts>]{<file>} -- include a figure from
% \datafigdir, degrading to a visible placeholder when the data directory has
% not been fetched. A missing artefact must not break the build: the manual has
% to compile from a bare clone.
\newcommand{\datafig}[2][width=\textwidth]{%
  \IfFileExists{\datafigdir #2}{%
    \includegraphics[#1]{#2}%
  }{%
    % \detokenize, not a bare \texttt{#2}: the filename was tokenized at the
    % call site, so its underscores carry catcode 8 and would be read as
    % subscripts. \detokenize re-emits them as printable catcode-12 characters.
    \fbox{\parbox{0.92\textwidth}{\centering\small\vspace{6pt}
      \textbf{Figure not available:} \texttt{\detokenize{#2}}\\[3pt]
      It lives in \texttt{\$LUPNT\_DATA\_PATH/docs/}, which this build could not
      find. See the figure notes in \texttt{latex\_docs/README.md} for how to
      regenerate it.\vspace{6pt}}}%
  }%
}
\usepackage[labelfont=bf,font=small]{caption}
\usepackage{enumitem}
\usepackage{microtype}

% Ragged-right X column: long identifiers in the spec tables must not be
% stretched into rivers.
\newcolumntype{L}{>{\raggedright\arraybackslash}X}

% Weighted ragged-right X column, used directly in a spec as `W{2}`: it takes
% twice the share of a plain `L`. The weights across one tabularx must sum to
% the number of X/L columns -- two columns can be W{1.4} W{0.6}, three can be
% W{1.5} W{1} W{0.5}, and so on. Reach for this instead of a fixed `l`/`r`/`c`
% column whenever a cell holds a phrase rather than a single token: a non-X
% column claims its content's natural width first, so one long cell there
% starves the real X columns until they wrap a word at a time.
\newcolumntype{W}[1]{>{\hsize=#1\hsize\raggedright\arraybackslash}X}
% Right-aligned weighted X, for a numeric comparison column that should still
% be able to wrap a trailing note.
\newcolumntype{Y}[1]{>{\hsize=#1\hsize\raggedleft\arraybackslash}X}

% A fixed-width ragged-right column, used as `P{2.1cm}`: for a token column
% (identifier, path) that is narrow on average but has an entry or two long
% enough that a bare `l` would widen the whole table.
\newcolumntype{P}[1]{>{\raggedright\arraybackslash}p{#1}}

% --- diagrams ---------------------------------------------------------------
\usepackage{tikz}
\usetikzlibrary{arrows.meta, positioning, fit, backgrounds, calc, shapes.geometric}

\tikzset{
  tsnode/.style   = {draw, rounded corners=2pt, minimum width=1.05cm,
                     minimum height=6mm, font=\small\sffamily, fill=white},
  boxnode/.style  = {draw, rounded corners=2pt, minimum width=2.4cm,
                     minimum height=8mm, font=\small\sffamily, fill=white,
                     align=center},
  constant/.style = {-, thick, black},
  table/.style    = {-, thick, black, dashed},
  linear/.style   = {-, thick, blue!65!black},
  model/.style    = {-, ultra thick, orange!85!black},
  integral/.style = {-, ultra thick, red!75!black, double},
  flow/.style     = {-{Latex[length=2mm]}, thick},
}

% --- admonitions ------------------------------------------------------------
% The reStructuredText sources use .. note:: / .. warning:: heavily; these
% reproduce them as light framed boxes without pulling in tcolorbox.
\usepackage{mdframed}
\usepackage{xcolor}

\definecolor{notebg}{RGB}{240,244,250}
\definecolor{noterule}{RGB}{ 70,110,170}
\definecolor{warnbg}{RGB}{253,244,238}
\definecolor{warnrule}{RGB}{200,110, 40}

\newmdenv[
  backgroundcolor=notebg, linecolor=noterule, linewidth=0.8pt,
  leftline=true, rightline=false, topline=false, bottomline=false,
  innerleftmargin=8pt, innerrightmargin=8pt,
  innertopmargin=6pt, innerbottommargin=6pt,
  skipabove=8pt, skipbelow=8pt
]{notebox}

\newmdenv[
  backgroundcolor=warnbg, linecolor=warnrule, linewidth=0.8pt,
  leftline=true, rightline=false, topline=false, bottomline=false,
  innerleftmargin=8pt, innerrightmargin=8pt,
  innertopmargin=6pt, innerbottommargin=6pt,
  skipabove=8pt, skipbelow=8pt
]{warnbox}

\newenvironment{lupntnote}{\begin{notebox}\noindent\textbf{Note.}\ \ignorespaces}
  {\end{notebox}}
\newenvironment{lupntwarning}{\begin{warnbox}\noindent\textbf{Warning.}\ \ignorespaces}
  {\end{warnbox}}

% --- source listings --------------------------------------------------------
\usepackage{listings}

\definecolor{codebg}{RGB}{248,248,248}
\definecolor{codekw}{RGB}{ 20, 70,160}
\definecolor{codestr}{RGB}{160, 60, 40}
\definecolor{codecomment}{RGB}{100,120,100}
\definecolor{codeframe}{RGB}{220,220,220}

\lstdefinestyle{lupnt}{
  backgroundcolor = \color{codebg},
  basicstyle      = \ttfamily\footnotesize,
  keywordstyle    = \color{codekw}\bfseries,
  stringstyle     = \color{codestr},
  commentstyle    = \color{codecomment}\itshape,
  numberstyle     = \tiny\color{gray},
  frame           = single,
  rulecolor       = \color{codeframe},
  framesep        = 4pt,
  xleftmargin     = 6pt,
  xrightmargin    = 0pt,
  breaklines      = true,
  breakatwhitespace = false,
  postbreak       = \mbox{\textcolor{gray}{$\hookrightarrow$}\space},
  columns         = fullflexible,
  keepspaces      = true,
  showstringspaces= false,
  captionpos      = b,
  aboveskip       = 8pt,
  belowskip       = 8pt,
  literate        = {->}{{->}}2 {<-}{{<-}}2 {~}{{\textasciitilde}}1
                    {_}{{\_}}1 {°}{{$^\circ$}}1 {±}{{$\pm$}}1
                    {μ}{{$\mu$}}1 {Δ}{{$\Delta$}}1 {→}{{$\rightarrow$}}1
                    {↔}{{$\leftrightarrow$}}1 {—}{{---}}1 {–}{{--}}1
                    {’}{{'}}1 {‘}{{`}}1 {“}{{``}}1 {”}{{''}}1,
}
\lstset{style=lupnt}

% Language aliases used by the chapters.  Each chapter selects one with
% \begin{lstlisting}[language=cpp] etc.; the names below are the only ones
% guaranteed to exist.
\lstdefinelanguage{cpp}[ANSI]{C++}{
  morekeywords={auto,constexpr,nullptr,override,using,namespace,template,
                typename,struct,class,enum,static,inline,const,noexcept,
                Real,Vec3,Vec6,VecX,MatX,Ptr,std},
}
\lstdefinelanguage{py}[]{Python}{
  morekeywords={None,True,False,self,async,await,match,case},
}
\lstdefinelanguage{yamlcfg}{
  keywords={true,false,null,name,type,agents,applications,measurements,
            simulation,dynamics,frame,epoch,duration,step,outputs},
  keywordstyle=\color{codekw}\bfseries,
  comment=[l]{\#},
  morestring=[b]',
  morestring=[b]",
  sensitive=false,
}
\lstdefinelanguage{shell}{
  keywords={cd,ls,git,cmake,pixi,pip,make,python,export,source,conda,sudo,apt,
            brew,doxygen,sphinx-build,tectonic},
  keywordstyle=\color{codekw}\bfseries,
  comment=[l]{\#},
  morestring=[b]",
  sensitive=true,
}
\lstdefinelanguage{bibtexlang}{
  keywords={@article,@inproceedings,@book,@techreport,@misc,@phdthesis},
  keywordstyle=\color{codekw}\bfseries,
  comment=[l]{\%},
  sensitive=false,
}

% --- inline literals --------------------------------------------------------
% \code   -- an identifier, type, enum value, CLI flag, or short expression
% \file   -- a path in the repository (breaks at / and _)
% \cpp    -- a C++ symbol (same rendering as \code; kept distinct so that a
%            future index pass can separate the two)
% \py     -- a Python symbol
\usepackage{url}
\urlstyle{tt}
% All four are robust so they survive unprotected in captions and section
% headings -- \path is fragile, and a bare \file{a_b.cc} in a \caption would
% otherwise fail with "Undefined control sequence".
\DeclareRobustCommand{\code}[1]{\texttt{#1}}
\DeclareRobustCommand{\cpp}[1]{\texttt{#1}}
\DeclareRobustCommand{\py}[1]{\texttt{#1}}
\DeclareRobustCommand{\file}[1]{\path{#1}}
\newcommand{\lupnt}{\textsc{LuPNT}}
\newcommand{\pylupnt}{\texttt{pylupnt}}

% Long typewriter identifiers that are allowed to wrap. `\code`/`\cpp` set a
% \texttt box that never breaks, so a 40-character API name pokes out of a
% table cell or into the margin. Two flavours, both underscore-safe:
%
%   \longid{Foo.barBaz_qux}  breaks only at dots, underscores and slashes --
%                            the tidy choice for a dotted or underscored name.
%   \brk{NONSPHERICALGRAVITY} breaks between any two characters -- the only
%                            option for a separator-less run of capitals that
%                            is itself wider than the column.
%
% Reach for either only when an identifier is too long to sit on one line;
% short ones stay \code.
\DeclareUrlCommand{\longid}{\urlstyle{tt}%
  \def\UrlBreaks{\do\/\do\_\do\.\do\-\do\:}}  % + C++ scope operator ::
\makeatletter
\DeclareUrlCommand{\brk}{\urlstyle{tt}%
  \def\UrlBreaks{\do\/\do\_\do\.\do\-%
    \do\A\do\B\do\C\do\D\do\E\do\F\do\G\do\H\do\I\do\J\do\K\do\L\do\M%
    \do\N\do\O\do\P\do\Q\do\R\do\S\do\T\do\U\do\V\do\W\do\X\do\Y\do\Z%
    \do\a\do\b\do\c\do\d\do\e\do\f\do\g\do\h\do\i\do\j\do\k\do\l\do\m%
    \do\n\do\o\do\p\do\q\do\r\do\s\do\t\do\u\do\v\do\w\do\x\do\y\do\z%
    \do0\do1\do2\do3\do4\do5\do6\do7\do8\do9}}
\makeatother

% Vectors and frequently used operators.  Chapters use \vec{} for 3-vectors
% and \mat{} for matrices so the whole manual stays typographically uniform.
\renewcommand{\vec}[1]{\bm{#1}}
\newcommand{\mat}[1]{\bm{#1}}
\newcommand{\uvec}[1]{\hat{\bm{#1}}}
\newcommand{\dd}{\mathrm{d}}
\newcommand{\transpose}{^{\mathsf{T}}}
\DeclareMathOperator{\diag}{diag}
\DeclareMathOperator{\atantwo}{atan2}
\DeclareMathOperator*{\argmin}{arg\,min}
\DeclareMathOperator{\Tr}{tr}
\DeclareMathOperator{\erf}{erf}

% Time scales typeset uniformly as upright abbreviations.
\newcommand{\TT}{\mathrm{TT}}
\newcommand{\TAI}{\mathrm{TAI}}
\newcommand{\TDB}{\mathrm{TDB}}
\newcommand{\TCB}{\mathrm{TCB}}
\newcommand{\TCG}{\mathrm{TCG}}
\newcommand{\TCL}{\mathrm{TCL}}
\newcommand{\LT}{\mathrm{LT}}
\newcommand{\UTC}{\mathrm{UTC}}
\newcommand{\UTone}{\mathrm{UT1}}
\newcommand{\GPST}{\mathrm{GPS}}

% --- bibliography and cross references --------------------------------------
\usepackage[numbers,sort&compress]{natbib}
\usepackage[hidelinks]{hyperref}
\usepackage{cleveref}

\crefname{chapter}{Chapter}{Chapters}
\crefname{section}{Section}{Sections}
\crefname{equation}{Eq.}{Eqs.}
\crefname{figure}{Figure}{Figures}
\crefname{table}{Table}{Tables}
\crefname{listing}{Listing}{Listings}

\setlength{\emergencystretch}{4em}  % absorb long unbreakable \texttt identifiers
